\section{Registro de configuración}

Los bloques siguientes son los comandos aplicados en cada router, tal como están en la carpeta \texttt{configs/} del repositorio.
Se ordenan por sistema autónomo.

\subsection*{Supuestos de implementación}

La topología entregada contiene los dispositivos con líneas dibujadas y sin cables reales, y cada router 1941 trae solo dos puertos Gigabit. El grupo la tomó como una referencia y supuso lo siguiente:
\begin{itemize}[nosep]
  \item Los enlaces se cablearon según la tabla siguiente. Los enlaces seriales usan un módulo HWIC-2T en la ranura 0 de cada router que lo necesita, con \cmd{clock rate 64000} en el extremo DCE (el router nombrado primero).
  \item Cada LAN se conecta a la interfaz \texttt{GigabitEthernet0/0} de su router mediante el switch; Base server se conecta al switch de Cypher para quedar en la misma LAN que Tripwire y Spycam.
  \item Los nombres de los equipos son los que aparecen en la topología entregada.
  \item Packet Tracer no soporta iBGP (al configurar un vecino del mismo AS responde ``Only external neighbors are supported''). Por eso, dentro del AS 40, Reyna, Omen y Sage corren OSPF (proceso 40) entre ellos y cada uno redistribuye entre OSPF y BGP con \cmd{redistribute bgp 40 subnets} y \cmd{redistribute ospf 40 match internal external 1 external 2}.
  \item En Raze y Killjoy la redistribución de OSPF hacia BGP se hace con \cmd{redistribute ospf 1 match internal}. Sin el filtro, Packet Tracer reexporta hacia el AS 20 las rutas externas aprendidas de Omega y se forma un bucle entre Neon y Raze (ver pregunta 3).
\end{itemize}

\begin{center}\footnotesize
\begin{tabular}{@{}llll@{}}
\toprule
Red & Cable & Extremo .1 & Extremo .2 \\
\midrule
1.0.0.0/8 & Cobre cruzado & Deadlock GigabitEthernet0/1 & Cypher GigabitEthernet0/1 \\
2.0.0.0/8 & Serial DCE & Deadlock Serial0/0/0 & Raze Serial0/0/0 \\
3.0.0.0/8 & Serial DCE & Cypher Serial0/0/0 & Killjoy Serial0/0/0 \\
4.0.0.0/8 & Serial DCE & Raze Serial0/0/1 & Killjoy Serial0/0/1 \\
5.0.0.0/8 & Cobre cruzado & Raze GigabitEthernet0/1 & Neon GigabitEthernet0/0 \\
6.0.0.0/8 & Cobre cruzado & Killjoy GigabitEthernet0/1 & Fade GigabitEthernet0/0 \\
7.0.0.0/8 & Cobre cruzado & Fade GigabitEthernet0/1 & Viper GigabitEthernet0/1 \\
8.0.0.0/8 & Cobre cruzado & Neon GigabitEthernet0/1 & Gekko GigabitEthernet0/1 \\
9.0.0.0/8 & Serial DCE & Viper Serial0/0/0 & Omen Serial0/0/0 \\
10.0.0.0/8 & Cobre cruzado & Omen GigabitEthernet0/1 & Sage GigabitEthernet0/0 \\
11.0.0.0/8 & Cobre cruzado & Reyna GigabitEthernet0/1 & Sage GigabitEthernet0/1 \\
12.0.0.0/8 & Serial DCE & Gekko Serial0/0/0 & Reyna Serial0/0/0 \\
13.0.0.0/8 & Serial DCE & Sage Serial0/0/0 & Brimstone Serial0/0/0 \\
\bottomrule
\end{tabular}

\end{center}

\newcommand{\registro}[2]{%
  \subsection{#1 (AS #2)}
  \lstinputlisting[caption={Configuración de #1}]{configs/#1.txt}}

\subsection*{AS 10 -- Base Alfa (OSPF área 0)}
\registro{Deadlock}{10}
\registro{Cypher}{10}
\registro{Raze}{10}
\registro{Killjoy}{10}

\subsection*{AS 20 -- Frontera Omega}
\registro{Neon}{20}
\registro{Fade}{20}

\subsection*{AS 30 -- Soporte táctico (con DHCP)}
\registro{Gekko}{30}
\registro{Viper}{30}

\subsection*{AS 40 -- Defensa y control}
\registro{Reyna}{40}
\registro{Omen}{40}
\registro{Sage}{40}

\subsection*{AS 50 -- Comando supremo}
\registro{Brimstone}{50}
